\chapter{Claim: Stock Market Returns Are Predictable}
\label{ch:08}

Every claim about forecasting, however simple or sophisticated the model behind it, has to pass the out-of-time test. This chapter builds the tools for it: loss functions and benchmarks, walk-forward designs that use only past information, the forms of leakage that let the future slip into a backtest, the statistics that decide whether one forecast beats another, and multiple-testing corrections for predictive model search. It then uses those tools on one of the most influential studies in empirical finance.

By the early 2000s, dozens of published papers had found variables that predicted the stock market's excess return: the dividend--price ratio, interest rates, yield spreads, inflation, corporate issuance, and more. Each had been significant in a regression on historical data, and predictability of the equity premium was widely taught as a fact. Would an investor who had used these variables in real time, knowing only what was known at each date, have forecast better than someone who simply assumed the average past return? (Amit Goyal and Ivo Welch, 2008)

The claim on trial is: \textbf{a predictor or model that is statistically significant in-sample will forecast the equity premium better than the historical mean out of sample, and well enough to be worth using.}

The chapter is built around Goyal and Welch (2008), who showed that popular predictors of the stock market forecast well in sample but poorly out of sample. It also draws on White's (2000) Reality Check for data snooping and on Campbell and Thompson (2008), who argued that small, sensible restrictions rescue some predictability. The lab in Section~\ref{sec:08-lab} replicates the Goyal--Welch horse race and White's test in order to check whether any predictor beats the historical average once we test it the way an investor would have used it.

By the end of the chapter you will be able to:
\begin{itemize}
  \item explain why in-sample $R^2$ flatters complex models, and why in-sample significance is a poor guide to out-of-sample performance;
  \item design a walk-forward evaluation using expanding or rolling windows with strict point-in-time information sets;
  \item recognize and prevent data leakage, applying purging and embargoing for overlapping targets and cross-validation;
  \item select natural hard-to-beat benchmarks and evaluate out-of-sample explanatory power ($R^2_{\text{OS}}$);
  \item execute hypothesis tests for forecast accuracy using Diebold--Mariano (for non-nested models) and Clark--West (for nested models);
  \item evaluate economic significance through certainty-equivalent return (CER) gains;
  \item apply White's (2000) Reality Check and Hansen's (2005) SPA test to evaluate data mining across multiple candidate forecasting models;
  \item analyze post-publication decay (McLean and Pontiff 2016) as a natural experiment;
  \item summarize the Goyal--Welch critique and evaluate the effectiveness of restricted and combined forecasts.
\end{itemize}

In the lab you will replicate Goyal and Welch's study with updated data, evaluating 14 classic predictors walk-forward from 1965 to the present, and apply White's Reality Check to a universe of predictive rules.

\section{Foundations: the basics of forecasting}
\label{sec:08-foundations}

\subsection{Information sets}
A forecast made at time $t$ for target $t+1$ must rely exclusively on the information set $\mathcal{F}_t$ available at date $t$: past prices, published macroeconomic releases, and models estimated on past data.

\subsection{Data splits: training, validation and test}
Forecasters divide data into three distinct chronological partitions:
\begin{itemize}
  \item \textbf{Training set}: used to estimate model parameters;
  \item \textbf{Validation set}: used to tune hyperparameters or select models;
  \item \textbf{Test set}: held out until final evaluation to measure performance out of sample.
\end{itemize}

\subsection{In-sample $R^2$ and its flattery}
For a model fitted to $T$ observations, the in-sample $R^2$ is:
\begin{equation}
  R^2 = 1 - \frac{\sum_{t=1}^T (y_t - \hat{y}_t)^2}{\sum_{t=1}^T (y_t - \bar{y})^2}
  \label{eq:08-r2}
\end{equation}
Because OLS explicitly minimizes the sum of squared residuals, in-sample $R^2$ can never decrease when candidate regressors are added. Consequently, in-sample $R^2$ flatters overfitted models, necessitating out-of-sample evaluation.

\section{In-sample fit versus out-of-sample performance}
\label{sec:08-insample}

The standard time-series predictive regression evaluates next-period excess return:
\begin{equation}
  r_{t+1} = \alpha + \beta x_t + \varepsilon_{t+1}
  \label{eq:08-pred}
\end{equation}
Goyal and Welch (2008) tested whether real-time recursive forecasts from Equation~\eqref{eq:08-pred} beat the simple historical mean $\bar{r}_t$. Most failed. Three primary forces open the gap between fit and forecast:
\begin{enumerate}
  \item \textbf{Estimation Error:} In finite samples with low signal-to-noise ratios, sampling noise in $\hat{\alpha}$ and $\hat{\beta}$ dominates the forecast.
  \item \textbf{Structural Instability:} Economic relationships evolve over time due to structural shifts, learning, and institutional changes.
  \item \textbf{Selection Bias:} Published predictors were selected because they fit historical data in-sample.
\end{enumerate}
In addition, Stambaugh (1999) bias inflates in-sample $t$-statistics when predictors $x_t$ are highly persistent ($\rho \approx 1$) and predictor shocks correlate with return shocks.

\section{Walk-forward evaluation: expanding and rolling windows}
\label{sec:08-walkforward}

A walk-forward evaluation replays history recursively:
\begin{enumerate}
  \item Select an initial forecast date $t_0$, establishing an initial estimation window.
  \item Estimate parameters using data available up to $t$.
  \item Forecast target $\hat{r}_{t+1}$ and record the benchmark forecast $\bar{r}_t$.
  \item Advance to $t+1$, update the training window, and repeat until the sample end.
\end{enumerate}

\begin{table}[htbp]
\centering
\small
\caption{Expanding vs. Rolling Window Architectures}
\label{tab:08-windows}
\begin{tabularx}{0.95\linewidth}{@{}>{\raggedright\arraybackslash}p{0.18\linewidth}>{\raggedright\arraybackslash}X>{\raggedright\arraybackslash}p{0.24\linewidth}>{\raggedright\arraybackslash}p{0.24\linewidth}@{}}
\toprule
Window Type & Training Data Used at Date $t$ & Key Advantage & Primary Limitation \\
\midrule
\textbf{Expanding} & All observations from start to $t$ & Uses maximum data; minimizes parameter estimation noise. & Adapts slowly if underlying data-generating process changes. \\
\textbf{Rolling} & Fixed window of most recent $w$ observations & Adapts quickly to structural shifts. & Subject to higher parameter estimation error; sensitive to choice of $w$. \\
\bottomrule
\end{tabularx}
\end{table}

Out-of-sample forecast accuracy is evaluated using the Campbell and Thompson (2008) out-of-sample $R^2$:
\begin{equation}
  R^2_{\text{OS}} = 1 - \frac{\sum_{t} \left( r_{t+1} - \hat{r}_{t+1} \right)^2}{\sum_{t} \left( r_{t+1} - \bar{r}_t \right)^2}
  \label{eq:08-r2os}
\end{equation}
$R^2_{\text{OS}} > 0$ indicates that model forecasts achieve lower mean squared error than the historical mean.

\subsection{Theoretical restrictions}
Campbell and Thompson (2008) showed that imposing economic restrictions improves out-of-sample performance:
\begin{itemize}
  \item Set $\hat{\beta} = 0$ if the estimated slope sign contradicts economic theory.
  \item Truncate negative equity premium forecasts at zero ($\hat{r}_{t+1} = \max(0, \hat{r}_{t+1})$).
\end{itemize}

\section{Data leakage: look-ahead, publication lags, purging and embargoing}
\label{sec:08-leakage}

Data leakage occurs whenever future information enters model training or specification:
\begin{itemize}
  \item \textbf{Full-Sample Estimation \& Preprocessing:} Estimating coefficients, standardizing features, or setting winsorizing cutoffs using full-sample statistics.
  \item \textbf{Hindsight Model Selection:} Choosing variables or hyperparameters based on full-sample metrics.
  \item \textbf{Publication Lags \& Revisions:} Ignoring release delays (e.g., CPI releases arrive weeks after month-end; macroeconomic data undergo revisions). Point-in-time archives (such as ALFRED) must be used.
  \item \textbf{Overlapping Targets \& Purging:} When predicting multi-period targets ($k$-period returns), training observations whose target windows overlap the forecast date contain future information. López de Prado (2018) defines \textbf{purging} as dropping these overlapping observations from the training set.
  \item \textbf{Cross-Validation \& Embargoing:} In $k$-fold cross-validation, training data follow test folds. \textbf{Embargoing} drops a block of training data immediately after the test fold to eliminate serial correlation leakage.
\end{itemize}

\section{Benchmarks that are hard to beat}
\label{sec:08-benchmarks}

\begin{table}[htbp]
\centering
\small
\caption{Natural Hard-to-Beat Forecasting Benchmarks}
\label{tab:08-benchmarks}
\begin{tabularx}{0.95\linewidth}{@{}>{\raggedright\arraybackslash}p{0.22\linewidth}>{\raggedright\arraybackslash}p{0.30\linewidth}>{\raggedright\arraybackslash}X@{}}
\toprule
Target Quantity & Natural Benchmark & Rationale \\
\midrule
Equity Premium & Historical mean return ($\bar{r}_t$) & Minimal estimation noise; avoids unstable parameter fits. \\
Exchange Rates / Prices & Random Walk ($y_{t+1} = y_t$) & Meese and Rogoff (1983) empirical result. \\
Interest Rates / Yields & Random Walk or Forward Rate & Yields are highly persistent. \\
Volatility & Realized Volatility / EWMA & Volatility exhibits strong clustering. \\
Macroeconomic Variables & Survey Consensus / Last Value & Professional surveys aggregate public information efficiently. \\
\bottomrule
\end{tabularx}
\end{table}

\section{Measuring forecast accuracy: Diebold--Mariano and Clark--West}
\label{sec:08-accuracy}

\subsection{Non-nested model comparison: the Diebold--Mariano test}
Diebold and Mariano (1995) test equal predictive accuracy between two non-nested forecasts based on loss differentials $d_t = e^2_{\text{bench},t} - e^2_{\text{model},t}$. The test statistic is a $t$-test on $\bar{d}$ using Newey--West HAC standard errors, with Harvey, Leybourne, and Newbold (1997) small-sample corrections.

\subsection{Nested model comparison: the Clark--West test}
When a predictive regression is compared against the historical mean, the models are nested. Under the null hypothesis ($\beta=0$), estimating $\beta$ introduces parameter estimation noise that elevates the model's mean squared error above the benchmark. The uncorrected Diebold--Mariano test is therefore undersized and biased against discovering true predictability.

Clark and West (2007) adjust for estimation noise by defining the adjusted loss differential:
\begin{equation}
  f_t = (r_{t+1} - \bar r_t)^2 - \left[(r_{t+1} - \hat r_{t+1})^2 - (\bar r_t - \hat r_{t+1})^2\right]
  \label{eq:08-cw}
\end{equation}
A one-sided $t$-test on the mean of $f_t$ tests whether the predictor carries significant predictive information.

\subsection{Economic value: the certainty-equivalent return (CER)}
For a mean-variance investor with risk aversion $\gamma$, portfolio weight in equities is $w_t = \frac{\hat{r}_{t+1}}{\gamma \hat{\sigma}^2_{t+1}}$. The gain in Certainty-Equivalent Return ($\Delta \text{CER}$) measures the risk-free return premium the investor would accept to replace benchmark-based allocation with model-based forecasts. Campbell and Thompson (2008) provide the approximation $\Delta E[R] \approx R^2_{\text{OS}} / S^2$, where $S$ is the benchmark Sharpe ratio.

\section{Multiple testing in forecasting: White's Reality Check}
\label{sec:08-reality}

When a researcher tests $N$ candidate forecasting rules or models, selecting the top performer induces data-mining bias. White (2000) introduced the \textbf{Reality Check for Data Mining} to test the null hypothesis that the best model among $N$ candidates does not outperform the benchmark:
\begin{equation}
  H_0: \max_{k=1, \dots, N} E[d_k] \le 0
  \label{eq:08-white-null}
\end{equation}
where $d_{k,t} = e^2_{\text{bench},t} - e^2_{k,t}$ represents the performance differential of model $k$ relative to the benchmark.

\subsection{The Reality Check protocol}
\begin{enumerate}
  \item Compute observed mean performance differentials $\bar{d}_k = \frac{1}{T} \sum_{t=1}^T d_{k,t}$ for all $k = 1, \dots, N$, and record $V^{\text{obs}} = \max_k \bar{d}_k$.
  \item Resample performance series jointly using the stationary block bootstrap (Politis and Romano 1994) across $B$ bootstrap iterations.
  \item Construct the null distribution by centering bootstrapped statistics:
  \begin{equation}
    V^{*,b} = \max_{k=1, \dots, N} \left( \bar{d}_{k}^{*,b} - \bar{d}_k \right)
  \end{equation}
  \item Calculate the $p$-value as the proportion of bootstrap iterations where $V^{*,b} \ge V^{\text{obs}}$.
\end{enumerate}
Hansen (2005) refined this test into the Superior Predictive Ability (SPA) test, which reduces sensitivity to poor candidate models.

\begin{callout}
\noindent \textbf{What Failed: A Century of Technical Trading Rules}

Brock, Lakonishok, and LeBaron (1992) documented significant profits from moving-average and trading-range rules applied to the Dow Jones Index from 1897 to 1986. Sullivan, Timmermann, and White (1999) re-evaluated these results across a universe of thousands of technical rules using White's Reality Check. In the original 1897--1986 sample the best rule still beat the benchmark after the correction for search. In the ten years that followed, 1987--1996, the best rule was no longer significant, and on S\&P 500 futures no rule beat the benchmark. The rules had worked in the sample where they were found but not after it, which is the pattern Chapter~\ref{ch:09} examines in detail.
\end{callout}

\section{Post-publication samples as natural experiments}
\label{sec:08-postpub}

McLean and Pontiff (2016) evaluated 97 published anomalies across three distinct periods:
\begin{enumerate}
  \item \textbf{In-Sample Period}: original discovery sample;
  \item \textbf{Post-Sample / Pre-Publication Period}: returns dropped by 26\%, reflecting data-mining bias;
  \item \textbf{Post-Publication Period}: returns dropped by an additional 32\% (58\% total decline), reflecting arbitrage activities by market participants.
\end{enumerate}

\begin{callout}
\noindent \textbf{What Failed: The Dividend Yield Out of Sample}

The dividend-price ratio exhibited strong in-sample predictive performance through the 1980s. During the 1990s bull market, however, dividend yields dropped to historical lows, generating persistent negative return forecasts while equity markets surged. Real-time forecasters relying on dividend yields would have remained underweight equities throughout one of the largest bull markets in history. Structural shifts in mean dividend payouts explain much of this out-of-sample breakdown (Lettau and Van Nieuwerburgh 2008; Goyal and Welch 2008).
\end{callout}

\section{Lab~\thechapter: Replicating Goyal and Welch (2008) and White's Reality Check}
\label{sec:08-lab}

\begin{callout}
\noindent \textbf{Lab~\thechapter{} at a glance.} \emph{Claim on trial:} In-sample significant predictors beat the historical mean out of sample. \emph{Data:} Goyal and Welch updated dataset (\texttt{PredictorData.xlsx}) containing 14 macro predictors (valuation ratios, interest rates, credit spreads, stock variance, net equity expansion, inflation). \emph{Code:} \texttt{ch05-out-of-sample/}.
\end{callout}

\subsection*{Lab setup and structure}
\begin{table}[htbp]
\centering
\small
\caption{The 14 Macroeconomic Predictors of Goyal and Welch}
\label{tab:08-goyal-predictors}
\begin{tabularx}{0.95\linewidth}{@{}>{\raggedright\arraybackslash}p{0.28\linewidth}>{\raggedright\arraybackslash}X@{}}
\toprule
Category & Predictors \\
\midrule
Valuation Ratios & Dividend--price ratio, dividend yield, earnings--price ratio, dividend payout ratio, book-to-market ratio. \\
Interest Rates & Treasury-bill rate, long-term government yield, long-term bond return, term spread. \\
Credit & Default yield spread (Baa minus Aaa), default return spread. \\
Other & Stock variance, net equity expansion, inflation (lagged 2 months). \\
\bottomrule
\end{tabularx}
\end{table}

\begin{itemize}
  \item \textbf{Part A: In-Sample Regressions.} Run univariate regressions for all 14 predictors; record slope, White $t$-statistic, and $R^2$.
  \item \textbf{Part B: Out-of-Sample Walk-Forward.} Generate recursive forecasts starting in 1965. Compute $R^2_{\text{OS}}$, Clark--West, and Diebold--Mariano statistics. Plot cumulative squared error differences.
  \item \textbf{Part C: Remedies.} Evaluate Campbell--Thompson sign restrictions, forecast combinations, and certainty-equivalent return gains ($\Delta \text{CER}$).
  \item \textbf{Part D: Post-Publication Natural Experiment.} Evaluate predictive performance in the post-2005 sample.
  \item \textbf{Part E: White's Reality Check.} Apply White's Reality Check to a simulated technical trading universe to measure search bias.
\end{itemize}

\section{Verdict}
\label{sec:08-verdict}

\begin{callout}
\noindent \textbf{Verdict: Fails}

For most single macroeconomic predictors, in-sample statistical significance fails to translate into out-of-sample forecast superiority over the historical mean benchmark (Goyal and Welch 2008; Goyal, Welch, and Zafirov 2024). Out-of-sample $R^2_{\text{OS}}$ metrics are predominantly negative, and success is concentrated in isolated historical shocks. Economic value survives weakly only when imposing sign restrictions and equal-weighted forecast combinations.
\end{callout}

\section{What investors can use}
\label{sec:08-investors}

\textbf{Rules of Thumb:}
\begin{itemize}
  \item Start from the historical mean: Use the historical mean return as the primary benchmark for asset returns.
  \item Prefer model combinations: Equal-weighted forecast combinations consistently beat individual predictive models by shrinking variance.
  \item Inspect cumulative error plots: Plot cumulative squared error differences to verify whether forecast gains are steady or concentrated in single shock events.
\end{itemize}

\textbf{6-Question Investor Audit Checklist:}
\begin{enumerate}
  \item Were forecasts generated strictly walk-forward without look-ahead bias?
  \item Were point-in-time publication delays and data revisions accounted for?
  \item Is the benchmark model appropriate and hard to beat?
  \item Was nested model comparison evaluated using the Clark--West test?
  \item Does performance survive multiple-testing corrections (White's Reality Check / SPA)?
  \item Does predictive efficacy persist in post-publication samples?
\end{enumerate}

\section*{Exercises}

\subsection*{Conceptual (all readers)}
\begin{enumerate}
  \item Explain why in-sample $R^2$ cannot decrease when adding variables, whereas out-of-sample $R^2_{\text{OS}}$ can be negative.
  \item A forecaster predicts 12-month returns monthly using past data. Explain which training observations must be purged to prevent leakage.
  \item Why is the uncorrected Diebold--Mariano test biased when comparing nested predictive models? What adjustment does Clark--West provide?
  \item Describe how White's Reality Check constructs a null distribution for the maximum performance statistic across $N$ models.
  \item Explain why an equal-weighted combination of 14 weak forecasts often beats every individual forecast out of sample.
\end{enumerate}

\subsection*{Lab extensions (technical readers)}
\begin{enumerate}
  \setcounter{enumi}{5}
  \item Replicate Part B of Lab~\thechapter{} across 14 macro predictors. Compute $R^2_{\text{OS}}$ and Clark--West $p$-values for 1965--2005 vs. post-2005 samples.
  \item Evaluate sensitivity to window choice by rerunning walk-forward forecasts using a 20-year rolling window versus an expanding window.
  \item Apply White's Reality Check to a universe of 500 technical trading rules applied to daily market index returns. Determine whether the top rule survives multiple-testing penalties.
\end{enumerate}

